\section{Work Plan and Schedule}

The research internship will last 12 months (September 2026 to August 2027) and will be divided into the following main phases:

\begin{enumerate}
    \item \textbf{Month 1-2: Integration and Study (Phase 1)} \\
    Settling at the host institution; familiarization with the \texttt{JuliaSmoothOptimizers} (JSO) infrastructure; study of trust-region methods and structured proximal operators.
    \item \textbf{Month 3-5: Hybrid Algorithmic Design (Phase 2)} \\
    Theoretical development of structured proximal operators as global preprocessors and their integration with the local coordinate descent and combinatorial refiners studied in Brazil. Preliminary feasibility analysis.
    \item \textbf{Month 6-8: Computational Implementation and Partial Report (Phase 3)} \\
    Parallel implementation of the designed algorithms in the Julia language, integrating the models developed in Brazil into the JSO framework. Elaboration of the annual FAPESP Scientific Report (February 2027).
    \item \textbf{Month 9-10: Benchmarking and Validation (Phase 4)} \\
    Execution of scalable test batteries on linear and logistic regression problems with $L_0$ penalization and group sparsity. Comparison against established packages in the literature.
    \item \textbf{Month 11-12: Consolidation and Dissemination (Phase 5)} \\
    Consolidation of the results, integration of the developed methodologies into the PhD thesis, and elaboration of the BEPE Final Scientific Report required by FAPESP.
\end{enumerate}

The knowledge transfer flow will take place through regular online follow-up meetings with the advisor in Brazil, ensuring that all code development and theoretical advances are directly incorporated into the PhD thesis.
